% !TeX root = ../main.tex
% ---------------------------------------------------------------------------
% Instrumentation, Monitoring and the Observational Method
% Scope extension: expands the monitoring section of Chapter 10 ( Chapter~\ref{ch:construction}).
% Key sources: Peck (1969) Rankine Chapter; Nicholson et al. (1999) CIRIA
% R185; Dunnicliff, Geotechnical Instrumentation for Monitoring Field
% Performance; EN 1997-1 clause on the observational method.
% ---------------------------------------------------------------------------
\chapter{Instrumentation, Monitoring and the Observational Method}
\label{ch:monitoring}


\section{Purpose and the observational method}
\label{sec:purpose-and-the-observational-method}

\subsection{Why monitor}
\label{sec:why-monitor}
\scopenote{Verification of design assumptions, construction control,
third-party protection, contractual record and knowledge for future projects.
Monitoring without a decision attached to it is waste.}

\subsection{The observational method}
\label{sec:observational-method}
\scopenote{Peck's formulation and the Eurocode 7 requirements: acceptable
limits of behaviour established in advance, a plan of contingency actions, and
the ability to implement them in time. Ab initio and best-way-out
applications.}

\section{Instrument types}
\label{sec:instrument-types}

\subsection{Deformation instruments}
\label{sec:deformation-instruments}
\scopenote{Convergence arrays and tape extensometers, borehole extensometers,
inclinometers, settlement cells and magnetic extensometers; ranges, precision
and installation requirements.}

\subsection{Stress, load and pressure instruments}
\label{sec:load-instruments}
\scopenote{Contact and embedment pressure cells, load cells on bolts and
anchors, strain gauges in shotcrete and segments, and the notorious difficulty
of measuring stress reliably.}

\subsection{Groundwater instruments}
\label{sec:piezometers}
\scopenote{Standpipe, vibrating wire and multi-level piezometers; response
time and its dependence on permeability; datum control.}

\section{Surface and remote monitoring}
\label{sec:surface-and-remote-monitoring}

\subsection{Surface and structural monitoring}
\label{sec:surface-monitoring}
\scopenote{Precise levelling, automated total stations and prism arrays,
tiltmeters, crackmeters and building survey; accuracy achievable in a live
urban environment.}

\subsection{Remote sensing and imaging}
\label{sec:remote-sensing}
\scopenote{InSAR for wide-area settlement, terrestrial LiDAR and
photogrammetry for profile and convergence, ground penetrating radar for
lining thickness and voids, computer vision for defect detection, and
microseismic sensing. Each is a data stream into the model of
Chapter~\ref{ch:digital-underground}, not an end in itself.}

\subsection{Virtual sensing and inferred quantities}
\label{sec:virtual-sensing}
\scopenote{Deriving quantities that are not directly measured --- lining
stress, leakage rate, damage extent --- by combining sparse instrumentation
with a calibrated model. Extends the reach of a monitoring array without
adding instruments.}

\section{Data acquisition and array design}
\label{sec:data-acquisition-and-array-design}

\subsection{Automated data acquisition}
\label{sec:data-acquisition}
\scopenote{Loggers, telemetry, real-time dashboards and alarm distribution;
data volume, power and survivability on a construction site.}

\subsection{Designing the monitoring array}
\label{sec:array-design}
\scopenote{Instrument selection driven by the design question, section spacing
along the drive, redundancy, installation timing relative to the face, and the
importance of capturing pre-construction baselines.}

\section{Trigger levels and interpretation}
\label{sec:trigger-levels-and-interpretation}

\subsection{Three outputs of a monitoring programme}
\label{sec:monitoring-outputs}
\scopenote{A programme that is working produces three things: an updated
diagnosis of current condition, a prediction of future performance and risk,
and a decision on whether, when and how to intervene. Data that produces none
of these is cost without benefit.}

\subsection{Trigger levels and contingency}
\label{sec:trigger-levels}
\scopenote{Green-amber-red thresholds derived from design predictions rather
than from convenience; the escalation and authority chain; and pre-agreed
contingency measures.}

\subsection{Data interpretation}
\label{sec:data-interpretation}
\scopenote{Separating signal from instrument drift, temperature effects and
construction disturbance; trend analysis; deciding when a reading indicates a
real change in behaviour.}

\section{Design updating and common failures}
\label{sec:design-updating-and-common-failures}

\subsection{Back-analysis and design updating}
\label{sec:monitoring-back-analysis}
\scopenote{Using measured response to update ground parameters and revise
support classes ahead of the face. Method in
Chapter~\ref{ch:advanced-numerical}.}

\subsection{Common failures of monitoring programmes}
\label{sec:monitoring-failures}
\scopenote{Instruments installed too late, no baseline, unread data, triggers
with no owner, and monitoring specified as a deliverable rather than as a
decision tool.}
